\documentclass[10pt]{article}
\usepackage[T1]{fontenc}
\usepackage{lmodern,amsmath,amssymb,amsthm,booktabs,microtype}
\usepackage[margin=0.8in]{geometry}
\usepackage[colorlinks=true,urlcolor=blue]{hyperref}
\theoremstyle{definition}
\newtheorem{theorem}{Theorem}
\newtheorem{lemma}{Lemma}
\newcommand{\cat}{\mathbin{\Vert}}
\title{Certified Least Values for OEIS A087354 at Indices 20--25}
\author{Qichao Wang\\{\small Hebei University of Technology, Tianjin, China}
\and Daoyu Dong\\{\small University of Electronic Science and Technology of China, Chengdu, China}}
\date{October 22, 2026}
\begin{document}
\maketitle
\begin{abstract}
OEIS A087354 asks for the least perfect square that is a decimal
concatenation of $n$ nonzero perfect-square blocks, allowing repetitions.
We determine the formerly unlisted value $a(20)$ and the values
$a(22),\ldots,a(25)$ using exact-integer, exhaustive branch-and-bound
certificates, and independently reproduce the reported value $a(21)$.
The proof combines a digit-length reduction with a complete prefix-square
interval test. This is an exact finite determination, \emph{not} an
all-index closed form or an analytic formula for general $a(n)$.
\end{abstract}

\section{New exact values}
Let $a(n)$ be the minimum positive square $S$ whose usual decimal
representation is a concatenation of exactly $n$ positive squares.
A one-digit nonzero square is in $D=\{1,4,9\}$, and a two-digit
square is in $E=\{16,25,36,49,64,81\}$.

\begin{theorem}\label{thm:values}
The following values are exact:
\begin{center}\small
\begin{tabular}{@{}rllr@{}}\toprule
$n$ & $a(n)$ & $\sqrt{a(n)}$ & digits\\\midrule
20 & \texttt{111125911149919941444} & \texttt{10541627538} &21\\
21 & \texttt{444411911999914911441} & \texttt{21081079479} &21\\
22 & \texttt{11199444494448114419449} & \texttt{105827427893} &23\\
23 & \texttt{144949449149149194448144} & \texttt{380722272988} &24\\
24 & \texttt{1111425414199449191191449} & \texttt{1054241629893} &25\\
25 & \texttt{41914991191911941194149136} & \texttt{6474178804444} &26\\\bottomrule
\end{tabular}
\end{center}
The value at $n=21$ was already noted on the OEIS page; the other
five values are new determinations relative to the cited OEIS entry
as accessed on October 22, 2026 \cite{oeis}.
\end{theorem}

\section{Digit-length reduction and exact pruning}
Every admissible $S$ has at least $n$ decimal digits because each square
block has at least one digit. Each value displayed in
Theorem~\ref{thm:values} has at most $n+1$ digits. Thus, to prove
minimality, it suffices to examine lengths $L=n$ and $L=n+1$
(and only $L=n$ for the 21-digit candidate at $n=21$).

For $L=n$, every block is in $D$. For $L=n+1$, precisely one block
is in $E$ and all remaining $n-1$ blocks are in $D$: more than one
extra digit or a block of three or more digits would exceed $n+1$.
Consequently the complete \emph{candidate languages} are
\[
 \mathcal W_{n,n}=D^n,\qquad
 \mathcal W_{n,n+1}=\bigcup_{j=0}^{n-1}D^j E D^{n-1-j}.
\]
Here each word means the decimal concatenation of its integer blocks.
The two finite languages describe \emph{all} possible candidates in
the relevant length ranges, including repeated blocks.

\begin{lemma}[Square-interval exclusion]\label{lem:interval}
Let $p$ be the integer represented by a prefix of $d$ digits of a
putative $L$-digit square, where $0\leq d\leq L$. Every decimal
completion of that prefix belongs to
\[
 I_{p,d,L}=[p10^{L-d},(p+1)10^{L-d}-1]\cap\mathbb Z.
\]
In particular, no completion is a square if
\[
 \left\lceil\sqrt{p10^{L-d}}\right\rceil
 > \left\lfloor\sqrt{(p+1)10^{L-d}-1}\right\rfloor.
\]
For searching up to a known admissible upper bound $U$, the same
assertion holds with the right endpoint replaced by
$\min\{(p+1)10^{L-d}-1,U\}$.
\end{lemma}
\begin{proof}
The interval is exactly the set of nonnegative integers extending the
prefix with $L-d$ decimal digits. An integer square exists in an
integer interval $[A,B]$ precisely if the least integer square root
at least $\sqrt A$ is no larger than the greatest integer square root
at most $\sqrt B$. Intersecting with $[0,U]$ does not discard any
candidate not exceeding $U$.
\end{proof}

\section{Reproducible exhaustive certificate}
For each $n$ listed, the verifier starts with prefix $p=0$, $d=0$
digits, and $k=0$ blocks, and extends only by allowed blocks in
$D$ and (if one extra digit remains) $E$. Before visiting the
children it applies Lemma~\ref{lem:interval}; an empty square-root
interval prunes the whole subtree. At $k=n$, a leaf is accepted only
if $d=L$ and $p$ is a perfect square. The integer upper bound $U$ is
set to the claimed value. A second branch-and-bound limit, the least
square already found, can only remove larger values. Thus the search
is exhaustive for all squares $\leq U$ at each admissible digit length.
The independent block-partition routine verifies the $n$-block
representation of each proposed value, so each upper bound is valid.
These facts prove Theorem~\ref{thm:values} by finite exhaustion.

In particular, for $n=20$ the first witness is
\[
\begin{split}
10541627538^2={}&1\cat1\cat1\cat1\cat25\cat9\cat1\cat1\cat1\cat4\\
 &\cat9\cat9\cat1\cat9\cat9\cat4\cat1\cat4\cat4\cat4.\end{split}
\]
There are no eligible 20-digit squares, and every 21-digit eligible
square smaller than this value is excluded. A complete per-length
search trace and the remaining factorizations are generated by
\texttt{verify\_minima.py} and \texttt{verify\_minima.cpp}.

The C++ implementation is an independent fast audit with native
128-bit integers; a floating-point square root is used \emph{only}
as a provisional estimate and is adjusted by exact integer-square
comparisons before any accept/reject decision. The Python version
uses \texttt{math.isqrt} throughout, without floating point. Both
implementations agree on all six values and all complete search
counts. Running the Python script with \texttt{--extended} includes
the more expensive $n=25$ case.

\paragraph{Scope.} This note does not prove any all-index formula for
$a(n)$. For each index treated, the bound $L\leq n+1$ follows from
an explicit witness; we do not assume it for arbitrary $n$.
It remains open here whether this bound, or a tractable exact
closed formula, holds at every index. The earlier separate
\href{https://github.com/dongxuelian2/oeis-papers/tree/main/papers/A087354}{existence manuscript}
proves that all indices are defined but not their minima.

\begin{thebibliography}{9}
\bibitem{oeis} OEIS Foundation Inc.,
\emph{A087354: Smallest square which is the concatenation of $n$
nonzero squares}, \url{https://oeis.org/A087354},
accessed October 22, 2026.
\end{thebibliography}
\end{document}
